\begin{definition}\label{def:collatz_graph}\lean{CollatzMapBasics.collatz_graph}\leanok
The Collatz graph $\mathcal{G}$ is a directed graph on $\mathbb{N}$ with an edge from $n$ to $T(n)$ for all $n \in \mathbb{N}$. We consider its associated simple undirected graph $\mathcal{G}'$.
\end{definition}

\begin{lemma}\label{lemma:collatz_graph_adj_iff}\lean{CollatzMapBasics.collatz_graph_adj_iff}\leanok
Two distinct vertices $a, b$ are adjacent in $\mathcal{G}'$ if and only if $b = T(a)$ or $a = T(b)$.
\end{lemma}
\begin{proof}\leanok
Immediate from the definition of the simple graph associated with a directed graph.
\end{proof}

\begin{lemma}\label{lemma:T_iter_reachable}\lean{CollatzMapBasics.T_iter_reachable}\leanok
For any $k, n \in \mathbb{N}$, $T^k(n)$ is reachable from $n$ in $\mathcal{G}'$.
\end{lemma}
\begin{proof}\leanok
By induction on $k$, following the edges $(T^i(n), T^{i+1}(n))$.
\end{proof}

\begin{lemma}\label{lemma:confluence_step}\lean{CollatzMapBasics.confluence_step}\leanok
If $T^i(a) = T^j(b)$ and $b$ is adjacent to $c$ in $\mathcal{G}'$, then there exist $i', j'$ such that $T^{i'}(a) = T^{j'}(c)$.
\end{lemma}
\begin{proof}\leanok
If $c = T(b)$, then $T^i(a) = T^j(b)$, so $T^{i+1}(a) = T(T^i(a)) = T(T^j(b)) = T^{j+1}(b) = T^j(c)$.
If $b = T(c)$, then $T^i(a) = T^j(b) = T^j(T(c)) = T^{j+1}(c)$.
\end{proof}

\begin{lemma}\label{lemma:confluence_of_reachable}\lean{CollatzMapBasics.confluence_of_reachable}\leanok
If $a$ and $b$ are reachable in $\mathcal{G}'$, then they have a common descendant under $T$, i.e., there exist $i, j$ such that $T^i(a) = T^j(b)$.
\end{lemma}
\begin{proof}\leanok
By induction on the path length between $a$ and $b$, using Lemma~\ref{lemma:confluence_step}.
\end{proof}

\begin{theorem}\label{theorem:collatz_graph_weakly_connected_iff_collatz}\lean{CollatzMapBasics.collatz_graph_weakly_connected_iff_collatz}\leanok
The Collatz graph $\mathcal{G}$ restricted to positive integers is weakly connected (i.e., any two $a, b \ge 1$ are reachable in $\mathcal{G}'$) if and only if the Collatz conjecture holds.
\end{theorem}
\begin{proof}\leanok
($\Rightarrow$) Assume connectivity. For any $n \ge 1$, $n$ is reachable from 1. By Lemma~\ref{lemma:confluence_of_reachable}, there exist $i, j$ such that $T^i(n) = T^j(1)$. Since 1 is in a cycle $\{1, 2\}$, $T^j(1) \in \{1, 2\}$, so $T^i(n)$ eventually hits 1. By Lemma~\ref{lemma:T_iter_implies_collatz_iter}, $n$ eventually hits 1 under $C$.
($\Leftarrow$) Assume the Collatz conjecture. For any $n \ge 1$, $n$ eventually hits 1 under $C$, and by Lemma~\ref{lemma:collatz_iter_to_T_iter}, it also hits 1 under $T$. By Lemma~\ref{lemma:T_iter_reachable}, $n$ is reachable from 1 in $\mathcal{G}'$. Thus any $a, b \ge 1$ are both reachable from 1, and hence reachable from each other.
\end{proof}
